% =============================================================================
% math_paper_template.tex
% Workshop: Scientific Writing with LaTeX, Overleaf & Prism
% Purpose:  A mathematics paper template with theorems, proofs, and advanced
%           equation environments.  Compiles without errors as-is.
% =============================================================================

\documentclass[12pt, a4paper]{article}

% --- Packages -----------------------------------------------------------------
\usepackage[utf8]{inputenc}
\usepackage[T1]{fontenc}
\usepackage{amsmath, amssymb, amsthm}   % Core math + theorem environments
\usepackage{mathtools}                   % Extensions to amsmath
\usepackage{graphicx}
\usepackage{booktabs}
\usepackage{hyperref}
\usepackage{geometry}
\geometry{margin=2.5cm}

% --- Theorem-like environments ------------------------------------------------
% amsthm provides \newtheorem and the proof environment.
% The first argument is the command name; the second is the printed label.
% [section] means numbering resets per section (e.g., Theorem 2.1).

\theoremstyle{plain}       % Used for theorem-type statements (italic body)
\newtheorem{theorem}{Theorem}[section]
\newtheorem{lemma}[theorem]{Lemma}          % Shares the theorem counter
\newtheorem{corollary}[theorem]{Corollary}
\newtheorem{proposition}[theorem]{Proposition}

\theoremstyle{definition}  % Upright body text
\newtheorem{definition}[theorem]{Definition}
\newtheorem{example}[theorem]{Example}

\theoremstyle{remark}      % Upright, lighter emphasis
\newtheorem{remark}[theorem]{Remark}

% --- Title block --------------------------------------------------------------
\title{On the Convergence of Iterative Methods\\
       for Nonlinear Operator Equations}
\author{%
    First A.\ Author\thanks{Department of Mathematics, University X.
    Email: \texttt{first@university-x.edu}}
    \and
    Second B.\ Author\thanks{Institute of Applied Mathematics, University Y.
    Email: \texttt{second@university-y.edu}}
}
\date{\today}

% ==============================================================================
\begin{document}
\maketitle

% --- Abstract -----------------------------------------------------------------
\begin{abstract}
    We study the convergence properties of a class of iterative methods
    applied to nonlinear operator equations in Banach spaces.  Under suitable
    Lipschitz continuity assumptions on the Fr\'{e}chet derivative, we
    establish both local and semi-local convergence results.  Our main theorem
    generalises classical results of Kantorovich and provides sharper error
    bounds.

    \medskip
    \noindent\textbf{Keywords:} nonlinear equations, iterative methods,
    Banach space, convergence analysis.

    \noindent\textbf{MSC 2020:} 65J15, 47H10, 47J25.
\end{abstract}

% ==============================================================================
\section{Introduction}
\label{sec:intro}

Let $X$ and $Y$ be Banach spaces and let $F \colon D \subseteq X \to Y$ be a
nonlinear operator.  We consider the problem of finding $x^{*} \in D$ such
that
%
\begin{equation}
    F(x^{*}) = 0.
    \label{eq:main_problem}
\end{equation}
%
Newton's method generates a sequence $\{x_n\}$ via
%
\begin{equation}
    x_{n+1} = x_n - F'(x_n)^{-1} F(x_n), \qquad n = 0, 1, 2, \dots
    \label{eq:newton}
\end{equation}
%
where $F'(x)$ denotes the Fr\'{e}chet derivative of $F$ at $x$.

Greek letters appear naturally: the Lipschitz constant $\gamma$, step size
$\alpha$, tolerance $\varepsilon > 0$, and domain radius $\delta$.

% ==============================================================================
\section{Preliminaries}
\label{sec:prelim}

\begin{definition}[Lipschitz Condition]
    \label{def:lipschitz}
    We say $F' \colon D \to \mathcal{L}(X, Y)$ satisfies a
    \emph{centre Lipschitz condition} at $x_0$ with constant $\omega > 0$ if
    %
    \begin{equation}
        \lVert F'(x) - F'(x_0) \rVert \leq \omega \lVert x - x_0 \rVert,
        \qquad \forall\, x \in D.
        \label{eq:lipschitz}
    \end{equation}
\end{definition}

\begin{definition}[$\omega$-Condition]
    Operator $F$ satisfies the \emph{$\omega$-condition} on $D$ if there
    exist constants $\beta, \eta, \omega > 0$ such that
    %
    \begin{align}
        \lVert F'(x_0)^{-1} \rVert           &\leq \beta,
            \label{eq:omega_a} \\
        \lVert F'(x_0)^{-1} F(x_0) \rVert    &\leq \eta,
            \label{eq:omega_b} \\
        \lVert F'(x) - F'(y) \rVert           &\leq \omega \lVert x - y \rVert,
            \quad \forall\, x, y \in D.
            \label{eq:omega_c}
    \end{align}
\end{definition}

% ==============================================================================
\section{Main Results}
\label{sec:main}

\begin{lemma}
    \label{lem:bound}
    Let $F$ satisfy the $\omega$-condition with parameters
    $(\beta, \eta, \omega)$ and set $h = \beta \omega \eta$.  If $h \leq
    \tfrac{1}{2}$, then for every $n \geq 0$,
    %
    \begin{equation}
        \lVert x_{n+1} - x_n \rVert
            \leq \frac{1}{2^{n}} \lVert x_1 - x_0 \rVert.
        \label{eq:bound}
    \end{equation}
\end{lemma}

\begin{proof}
    We proceed by induction on $n$.  The base case $n = 0$ is immediate from
    the definition of $\eta$.  Assume the bound holds for some $n \geq 0$.
    By the Newton step~\eqref{eq:newton} and
    estimate~\eqref{eq:lipschitz}, we obtain
    %
    \begin{align}
        \lVert x_{n+2} - x_{n+1} \rVert
            &= \lVert F'(x_{n+1})^{-1} F(x_{n+1}) \rVert \notag \\
            &\leq \frac{\omega}{2}
                  \lVert x_{n+1} - x_n \rVert^{2}
                  \cdot \lVert F'(x_{n+1})^{-1} \rVert \notag \\
            &\leq \frac{1}{2} \lVert x_{n+1} - x_n \rVert,
    \end{align}
    %
    which completes the induction.
\end{proof}

\begin{theorem}[Semi-local Convergence]
    \label{thm:convergence}
    Under the hypotheses of Lemma~\ref{lem:bound}, the sequence $\{x_n\}$
    defined by~\eqref{eq:newton} converges to a solution $x^{*}$ of
    $F(x^{*}) = 0$.  Moreover, the error satisfies
    %
    \begin{equation}
        \lVert x_n - x^{*} \rVert
            \leq \frac{2\eta}{2^{n} - 1},
            \qquad n \geq 1.
        \label{eq:error}
    \end{equation}
\end{theorem}

\begin{proof}
    Since $\{x_n\}$ is Cauchy by Lemma~\ref{lem:bound} and $X$ is complete,
    the sequence converges to some $x^{*} \in \overline{B}(x_0, r)$ where
    %
    \[
        r = \frac{1 - \sqrt{1 - 2h}}{\omega \beta}.
    \]
    %
    Continuity of $F$ then gives $F(x^{*}) = 0$.  The error
    bound~\eqref{eq:error} follows by summing the geometric series from
    Lemma~\ref{lem:bound}.
\end{proof}

\begin{corollary}
    \label{cor:quadratic}
    If $F'$ is Lipschitz continuous on $D$ and $h < \tfrac{1}{2}$, then
    Newton's method converges \emph{quadratically}:
    %
    \begin{equation}
        \lVert x_{n+1} - x^{*} \rVert
            \leq C \lVert x_n - x^{*} \rVert^{2},
    \end{equation}
    %
    where $C = \tfrac{\omega \beta}{2}$.
\end{corollary}

% --- Matrices -----------------------------------------------------------------
\subsection{Matrix Notation Example}
\label{subsec:matrix}

The Jacobian of a system $\mathbf{F} \colon \mathbb{R}^{n} \to
\mathbb{R}^{n}$ can be written as

\begin{equation}
    J(\mathbf{x}) =
    \begin{pmatrix}
        \dfrac{\partial F_1}{\partial x_1} & \cdots &
            \dfrac{\partial F_1}{\partial x_n} \\[8pt]
        \vdots & \ddots & \vdots \\[4pt]
        \dfrac{\partial F_n}{\partial x_1} & \cdots &
            \dfrac{\partial F_n}{\partial x_n}
    \end{pmatrix}.
    \label{eq:jacobian}
\end{equation}

% --- Various math demonstrations ----------------------------------------------
\subsection{Additional Notation}
\label{subsec:notation}

Fractions and subscripts:
\[
    a_{n} = \frac{n!}{\Gamma(n + \alpha + 1)}, \qquad
    S_{N} = \sum_{k=0}^{N} \frac{(-1)^{k}}{(2k+1)!} \, x^{2k+1}.
\]

Operators and Greek letters:
\[
    \lim_{n \to \infty} \left(1 + \frac{1}{n}\right)^{n} = e, \qquad
    \int_{0}^{\pi} \sin\theta \, d\theta = 2.
\]

% ==============================================================================
\section{Conclusion}
\label{sec:conclusion}

We have established semi-local convergence of Newton's method under
a centre Lipschitz condition, obtaining error bounds that improve upon the
classical Kantorovich theory.  Future work will explore higher-order methods
and relaxed smoothness assumptions.

% --- Bibliography (self-contained using thebibliography) ----------------------
\begin{thebibliography}{9}

\bibitem{kantorovich1948}
    L.~V.~Kantorovich,
    \emph{On Newton's method for functional equations},
    Doklady Akad.\ Nauk SSSR, \textbf{59} (1948), 1237--1240.

\bibitem{ortega1970}
    J.~M.~Ortega and W.~C.~Rheinboldt,
    \emph{Iterative Solution of Nonlinear Equations in Several Variables},
    Academic Press, New York, 1970.

\bibitem{argyros2008}
    I.~K.~Argyros,
    \emph{Convergence and Applications of Newton-type Iterations},
    Springer, New York, 2008.

\bibitem{deuflhard2011}
    P.~Deuflhard,
    \emph{Newton Methods for Nonlinear Problems},
    Springer Series in Computational Mathematics, vol.~35, 2011.

\end{thebibliography}

\end{document}
